\chapter{Embodiment}\index{Embodiment}
\printchapterversion{04_Embodiment}

\newthought{Embodied AI and the Physical Contract}

\marginnote{%
    \begin{flushright}
    \newthought{Über das Marionettentheater}\\
    Heinrich von Kleist (1810)\\
    {\color{black}\qrcode[height=0.9cm]{https://de.wikisource.org/wiki/Ueber_das_Marionettentheater}}
    \end{flushright}%
}


% As a leading group, you are required to moderate this session. Add margin notes in this script in case you want to add any additional questions or points to the discussion.

\newthought{A dancer argues} that marionettes possess more grace than trained human performers,
because they lack the self-consciousness that disrupts human movement.\\

\newthought{Embodied AI} in physical space (robots, prosthetics, surgical systems) raises
expectations and moral responsibilities that purely software agents do not.

\vspace{3.5em}

\begin{tabular}{p{3.8cm} p{6.5cm}}
    \textbf{Required reading} &
        \textit{Über das Marionettentheater}\index{Kleist, Heinrich von}\index{Uber das Marionettentheater@\"{U}ber das Marionettentheater}\cite{kleist1810marionettentheater}
        by Heinrich von Kleist (1810; Reclam UB~9905) \\[3pt]
    \textbf{Preparation}      &
        Complete the Guided Reading Sheet individually before the session. \\
\end{tabular}

\vspace{3.5em}

\section*{Guided Reading Sheet}

\noindent Read \textit{Über das Marionettentheater} before the session. Work through the
questions below individually, then bring your notes to the discussion.

\begin{enumerate}
    \item The dancer argues that marionettes possess a grace that
    trained human performers cannot achieve, because they have no
    self-consciousness to disturb their movement. What does this claim
    imply about the relationship between embodiment and intelligence?
    Does having a body always aid or complicate performance?

    \item Moravec's paradox observes that tasks easy for humans
    (catching a ball, navigating a crowded room) are harder for
    machines than abstract reasoning like chess. What does this suggest
    about the relationship between intelligence and the body?

    \item A care robot that assists elderly patients shares physical
    space with them, touches them, and responds to their distress. Does
    its physical presence create obligations that a voice assistant
    does not have? If so, where do those obligations come from? One
    way to approach the question is by analogy to animal ethics, where
    we already have a framework for embodied, non-verbal, non-human
    beings: Singer\cite{singer1975animal} argues that capacity for
    suffering, not rationality or language, grounds moral
    consideration. When a robot can be harmed, distressed, or worn
    down, animal ethics becomes the nearest map. Similar lines have
    been considered by Asimov\cite{asimov1951foundation} and more
    recently by Gunkel\cite{gunkel2018robot}.

    \item Consider the concept of affordances: the actions an
    environment or object makes available to an agent. A robot's
    affordances include physical force and spatial presence; software's
    do not. How do the affordances of a robot differ from those of
    software, and what are the ethical implications of those
    differences?

    \item Kleist suggests that perfect grace requires either zero
    consciousness (a puppet) or infinite consciousness (a god), not
    the reflective, self-interrupting awareness humans have. Where
    would an embodied AI sit on this spectrum? Could a robot achieve
    Kleistian grace that humans cannot, and what would that mean for
    how we regard it morally?\\
\end{enumerate}



% ============================================================
% Running case options for Embodiment (draft, choose one later).
% Placed below the Guided Reading Sheet per convention. Uncomment
% the chosen block, delete the rest. Confirm dates/numbers first.
% ------------------------------------------------------------
% OPTION A, physical force and the blamed operator (closest to the Flash Crash):
% \vspace{3.5em}
% \noindent\textbf{Running case: the Uber test car.}
% \begin{itemize}
%     \item 18 March 2018, Tempe, Arizona: a self-driving Uber strikes
%           and kills Elaine Herzberg as she wheels a bicycle across a road
%           at night, the first pedestrian killed by an autonomous vehicle.
%     \item The system saw her seconds before impact but, tuned to ignore
%           false alarms, did not brake. The safety driver, meant to watch
%           the road, was looking away.
%     \item Uber settled with the family within days and was not charged.
%           The safety driver was charged with negligent homicide and later
%           pleaded guilty to endangerment.
% \end{itemize}
% ------------------------------------------------------------
% OPTION B, industrial death, no one intended it (German setting):
% \vspace{3.5em}
% \noindent\textbf{Running case: the Baunatal robot.}
% \begin{itemize}
%     \item July 2015, Volkswagen plant, Baunatal: a 22-year-old contractor
%           setting up a stationary assembly robot is seized and crushed
%           against a metal plate, and dies of his injuries.
%     \item The robot did exactly what it was programmed to do. The man was
%           inside its safety cage to configure it, so nothing stopped it.
%     \item Volkswagen called it human error. Prosecutors weighed whether
%           anyone could be charged for a death no one intended.
% \end{itemize}
% ------------------------------------------------------------
% OPTION C, machine mortality, feeds the Backup Ransom directly:
% \vspace{3.5em}
% \noindent\textbf{Running case: the AIBO funerals.}
% \begin{itemize}
%     \item Sony sold the AIBO robot dog from 1999, then ended repair
%           support in 2014, leaving thousands of owners with machines that
%           could no longer be fixed.
%     \item At Kofukuji temple in Chiba, Buddhist priests hold funerals for
%           dead AIBOs, chanting for machines their owners grieved as family.
%     \item Working parts are harvested to keep other AIBOs alive. Owners
%           speak of organ donation, not recycling.
% \end{itemize}
% ------------------------------------------------------------
% OPTION D, care and medicine, diffuse responsibility:
% \vspace{3.5em}
% \noindent\textbf{Running case: the surgical robot.}
% \begin{itemize}
%     \item The da Vinci system lets a surgeon operate through robotic arms
%           from a console, now used in millions of procedures worldwide.
%     \item A study of US regulator data linked 144 deaths and 1{,}391
%           injuries to robotic surgery between 2000 and 2013: device faults,
%           electrical arcs, and pieces left inside patients.
%     \item When an arm fails mid-operation, responsibility is split across
%           surgeon, hospital, and manufacturer, and often no single party
%           is clearly at fault.
% \end{itemize}
% ============================================================

\section*{In-Class Experiment: The Backup Ransom}

\noindent A short escalation to anchor the debate in shared experience. Run
before the discussion. The class faces the same deletion four times, each
with one more property of individuality in place, and marks after each
round: \textit{if we delete or destroy it here, has an individual died?
Rate~0 to~10.}

\begin{itemize}
    \item \textbf{Rung 1, perfect backup.} Pure software, with a current
    and identical backup that boots tomorrow. Deleting it costs nothing
    that is not restored. This is the baseline; almost no one calls it a
    death.
    \item \textbf{Rung 2, stale backup.} The snapshot is weeks old and the
    system has learned since. Restoring it returns an earlier self, and what
    is lost is everything it has become in between. Does a continuous
    history make the loss real?
    \item \textbf{Rung 3, no backup, still software.} A unique learned
    state on a server, unrecoverable, but with no body. Is irreversibility
    alone enough to individuate, without a body?
    \item \textbf{Rung 4, embodied and uncopyable.} The only instance,
    running in one robot body in volatile memory, physically impossible to
    copy. Destroy it and it is gone by necessity, not by policy.
\end{itemize}

\newthought{The sharp variant} sits between Rung~3 and Rung~4 and divides
opinion most: the system is embodied, but its mind-state is restorable into
an identical new body. Destroy the body, reload the same weights into an
identical chassis. Death, or repair? Whoever answers death holds that the
body individuates; whoever answers repair holds that the mind-state does.

\newthought{Debrief.} Put the four numbers on the board and compare where
students crossed from ``restore it'' to ``it died''. Cross at Rung~2 and
you hold that continuous memory makes the self; at Rung~3, irreversibility;
at Rung~4, an uncopyable body; never, and you deny there was any individual
to lose. The disagreement is the point. Individuality is not one property
but several (memory, irreversibility, an uncopyable body) that ordinarily
never come apart. Software separates them; a body holds them together.

\section*{What a Body Makes}

\noindent The experiment isolates the first of two properties a body
confers. Call it individuality: the word individual means undividable, and
only something undividable can be lost as a whole. A running process is
endlessly divisible, and nothing about it is ever finally lost; a body
cannot be divided or restored, and so it can die.

\newthought{The second property} is a surface that can be hurt. A capacity
to suffer requires something that can be damaged and not reloaded from
backup, and a body supplies it. Animal ethics reached this pivot two
centuries ago: Bentham asked not whether an animal can reason but whether
it can suffer, and Singer built moral standing on that capacity alone.
Kleist's marionette is the counterexample: a body of perfect grace with no
self inside it, embodiment without a subject.

\newthought{The question the chapter forces} is whether giving an AI a body
changes its interface or its kind. If a body makes an individual, and a
vulnerable body makes a sufferer, then embodiment is not a feature added to
a tool. It may be the step that turns a tool into someone we can wrong.

\section*{The Body as a Bottleneck}

\marginnote{\newthought{Transhumanism.} The view that the human body and
mind are a starting point to be engineered past, through augmentation,
life extension, and eventually the merger of human and machine.}

\noindent This chapter has argued that a body is what makes an individual
and a sufferer. Transhumanism argues the opposite: the body is a
bottleneck, a mortal and error-prone substrate to be upgraded and
eventually left behind. Where this chapter asks what we owe an embodied
machine, transhumanism asks how quickly we can stop being merely embodied
ourselves.

\marginnote{\newthought{Technological singularity.} A hypothesised point at
which machine intelligence, improving itself, outruns human understanding,
after which the future becomes unpredictable from where we now stand.}

\newthought{Ray Kurzweil} gives the optimistic case. His law of
accelerating returns reads technological progress as exponential rather
than linear, and extrapolates it to a singularity around 2045: the moment
machine intelligence exceeds our own and we merge with it, transcending
biology rather than being replaced by it.\cite{kurzweil2005singularity}
Death becomes an engineering problem, and the vulnerable body the
experiment treated as the seat of individuality becomes optional.

\newthought{Peter Thiel} shares the aspiration and distrusts the curve. He
funds life extension and treats death as a problem to be solved rather than
accepted, yet argues that no exponential arrives on its own: the future is
built by people who choose to build it, not delivered by a
graph.\cite{thiel2014zero} The singularity, on this reading, is not a
prophecy to wait out but a project that can also fail, and the people
pursuing it are as embodied and mortal as anyone the chapter has discussed.

\marginnote{\newthought{Individualism.} The tradition that treats the
single, unrepeatable person as the basic unit of moral and political worth.
A self that can be copied or merged offers no single instance for that
worth to attach to.}

\newthought{What makes a human human?} The chapter asked what makes an
individual and answered with the undividable, vulnerable body.
Transhumanism forces the sharper form of the question. Kurzweil's merger
makes a mind substrate-independent, and a substrate-independent mind is
copyable and restorable, exactly the properties the Backup Ransom said an
individual lacks. Lift the self off the body and you lift off the
individuality the body guaranteed: the tradition of individualism, which
treats the single unrepeatable person as the seat of moral worth, loses the
single instance it was built on. Uploading need not preserve the person; it
may only produce a copy that believes it is you, while the original still
dies. The experiment we ran on a robot runs just as well on us.

\newthought{The tension is the point.} If the body is a bottleneck, the
right response to an ageing, breakable human is to engineer past it. If the
body is what makes an individual who can be wronged, then engineering it
away is not liberation but the quiet removal of the very thing that made us
someone. The debate below is where that disagreement gets its hearing.

\section*{Opinion Landscape and Debate}

\noindent \textbf{Format:} Surrounded-style debate. Follow the shared
debate protocol for the speaker's list rules and moderator scripts.

\medskip
\noindent \textbf{Opening question:} Does having a physical body change
what we owe a machine, or what a machine owes us?
\newthought{Opinion~A, Equivalence.}\quad Embodied AI operating in
intimate or high-risk physical settings must meet the same professional
liability standards as the human practitioners they are deployed alongside.
\vspace{1.5em}


\medskip
\noindent\textbf{The Patient Advocate} argues that deploying a robot in intimate
    care settings requires a higher standard of informed consent.
    \begin{itemize}
        \item \textit{A care relationship involves physical proximity, trust,
        and vulnerability. My client never gave explicit, informed consent to
        physical handling by a non-human agent. That consent cannot be
        implied from a general admission form.}
        \item \textit{The power asymmetry between a frail elderly patient
        and a care institution is significant. Consent obtained in that
        context cannot be considered freely given without specific disclosure
        about robotic care.}
        \item \textit{We are not arguing that robots cannot be used in care.
        We are arguing that the standard of consent for their use in intimate
        physical tasks must be at least as rigorous as it is for medical
        procedures.}
    \end{itemize}

\medskip
\noindent\textbf{The Ethicist} questions whether physical autonomy in care should
    ever be delegated to a non-human agent, regardless of performance metrics.
    \begin{itemize}
        \item \textit{We are focused on liability, but the prior question is
        whether this deployment was ethical to begin with. Why did we decide
        that intimate physical care (which requires responsiveness to
        pain, distress, and dignity) is an appropriate domain for
        automation?}
        \item \textit{Robots are being deployed in care because human carers
        are undervalued and in short supply. The liability debate is a
        distraction from that structural choice. We are managing the
        consequences of a political economy decision by holding the wrong
        parties accountable.}
        \item \textit{Whatever liability framework we adopt here will create
        incentives for future deployment decisions. I want to ask: what does
        the framework we design signal about what we value: efficiency,
        or dignity?}
    \end{itemize}



\newthought{Opinion~B, Tool.}\quad An AI is a tool; the relevant
liability flows through the humans and institutions that deploy it,
not through the system itself.
\vspace{1.5em}


\medskip
\noindent\textbf{The Manufacturer} argues the robot performed within specification;
    the care home misconfigured its environment.
    \begin{itemize}
        \item \textit{Our robot was certified to the highest internationally
        recognised safety standards for collaborative robotics. The care home
        deployed it in a space that violated the operational envelope
        specified in the manual. That is a deployment failure, not a
        manufacturing failure.}
        \item \textit{Care home staff modified the robot's proximity sensor
        settings without authorisation. That modification voids the warranty
        and transfers liability entirely.}
        \item \textit{If manufacturers are held liable for every way a
        complex system can be misused, no company will develop assistive
        robotics. The regulatory consequence of your position is that the
        technology disappears.}
    \end{itemize}

\medskip
\noindent\textbf{The Care Home Director} argues the manufacturer overstated
    the robot's capability for complex physical tasks.
    \begin{itemize}
        \item \textit{The manufacturer's marketing material (the brochure,
        the demonstration video, the sales pitch) showed exactly this task
        being performed in a facility like ours. We deployed the product as
        marketed.}
        \item \textit{There is a systematic gap between tested performance
        in controlled environments and real-world performance in a working
        care facility. The manufacturer knows this gap exists. They chose
        not to disclose it.}
        \item \textit{We are responsible for the welfare of our residents.
        We relied on the manufacturer's assurances in good faith. If those
        assurances were misleading, the responsibility lies with the party
        that made them.}
    \end{itemize}


\newthought{Key Learnings}
Embodiment is not merely a technical property but a social one: a physical agent occupies space,
exerts force, and triggers moral and legal responses that disembodied software does not. Moravec's
paradox reminds us that our intuitions about what is ``easy'' or ``hard'' for intelligence are shaped
by our embodied experience, and are regularly wrong. Kleist's dialogue makes visible the question
that liability law must eventually answer: what work is the concept of ``human'' doing in law, and
should physical presence, force, and risk be sufficient to extend it? The deployment of robots in
care, policing, and domestic settings is a present policy challenge; the frameworks need to be
built now, not after the first serious incident.

\medskip
\noindent\textit{Teaser: when a machine has a body, acts with force, and
answers to no one, what do we owe it?}

\clearpage
